\section{Random-design growth, fallback, and expected length}
\label{sec:design-stability}

Design failure does not itself invalidate coverage: the source-failure branch
uses the protected target band, and target-design failure returns the full
parameter interval. Conditional outcome guarantees and the unconditional
composition event can therefore be combined for every $K$. Conditions making
all designs good are needed for useful borrowing and length rates, rather
than for removing a fallback error from the coverage bound.

\subsection{An explicit sufficient design condition}

Let $d=p+1$, assume independent patients within each site, and suppose uniformly
\[
 \|\phi(X)\|_2^2\le L_d,\qquad
 \lambda_{\min}\{\E_j[\ind(A=a)\phi(X)\phi(X)^T]\}\ge\kappa>0.
\]
For $n_j$ patients, matrix Chernoff \citep[Corollary 5.2]{tropp2012} gives
\[
 \Prb\{\lambda_{\min}(X_{ja}^TX_{ja})<n_j\kappa/2\}
 \le d\exp\{-n_j\kappa/(8L_d)\}.
\]
The bound concerns the total site sample, so no deterministic arm sample size
is assumed. A union bound over both arms at the target and all sources gives
\[
 q_{\rm lower}\le 2(K+1)d\exp\{-n_{\min}\kappa/(8L_d)\}.
\]
On this event's complement,
\[
 \max_i h_{jai}\le\frac{2L_d}{n_j\kappa},\qquad
 \operatorname{cond}(X_{ja})\le\sqrt{2L_d/\kappa}.
\]
The first inequality ensures the leverage guard when its right-hand side is
below $.98$; it also rules out a saturated arm fit. To control the fixed
condition-number guard as dimension grows, additionally assume the population
arm-Gram maximum eigenvalues are at most $\Lambda$, uniformly. The upper
Chernoff bound gives, for $c_+=2\log2-1$,
\[
 q_{\rm upper}\le 2(K+1)d\exp\{-c_+n_{\min}\Lambda/L_d\}.
\]
On the combined event, $\operatorname{cond}(X_{ja})\le2\sqrt{\Lambda/\kappa}$,
which must lie below the numerical guard. Since $\Lambda\ge\kappa$,
\[
 q_{n,p,K}:=q_{\rm lower}+q_{\rm upper}
 \le4(K+1)d\exp\{-n_{\min}\kappa/(8L_d)\}.
\]
These statistical bounds concern exact arithmetic. The implementation separately
checks balance residuals and records any numerical failure as a fallback.

For bounded coordinates, $L_d=d$ is available. Thus
$d\log(Kd)/(n_{\min}\kappa)\to0$ is one sufficient joint-growth condition.
It is deliberately conservative, and its probability bound can be vacuous at
finite sample sizes where the observed designs work well. The experiments
report the sufficient bound and empirical failures separately.

Let $B_t=\|b_t\|_2$. On good designs and with leverage bounded by $h_*<1$,
\begin{align*}
 S_{ja}&\le\frac{2B_t^2}{n_j\kappa},&
 \|\gamma_{ja}\|_\infty&\le\frac{2\sqrt{L_d}B_t}{n_j\kappa},\\
 G_{ja},F_{ja}&=O\!\left(\frac{L_dB_t^4}{n_j^3\kappa^3}\right),&
 \max_i d_{jai}&=O\!\left(\frac{L_dB_t^2}{n_j^2\kappa^2}\right).
\end{align*}
Here $F_j\le G_j$, and
$G_j\le\|\gamma_j\|_\infty^2\|\gamma_j\|_2^2/(1-h_*)^2$.
For the common target mean, a bounded-difference argument supplies
\[
 \|b_t-\E_t\phi(X)\|_2
 \le \sqrt{L_d/n_t}+\sqrt{2L_d\log(1/\delta_t)/n_t}
\]
with probability at least $1-\delta_t$. This target event is paid for once,
not once per source. Bounded population feature-mean norm and
$L_d\log(1/\delta_t)/n_t=O(1)$ give $B_t=O(1)$.

\subsection{From conditional rates to expected length}

Under exact linear means, the zero-diagonal representation also gives
\[
 \Var(\widehat D\mid\mathcal D)
 \le\|Qm\|_2^2+\|Q\|_F^2/8\le LD+Q_0.
\]
The linear and off-diagonal quadratic parts have zero covariance for independent
centered outcomes, and bounded outcome variances are at most $1/4$.
At $D=0$, $\E|\widehat D|\le\sqrt{Q_0}$. Put $x=\log(1/\alpha_D)$ and
$M=\sqrt{Q_0}+cx$. The inversion formula implies
\[
 \E(U_D\mid\mathcal D)
 \le M+xL+\sqrt{2xLM+x^2L^2+2xQ_0}.
\]
This follows from $|\widehat D+cx|$ and Jensen's inequality; an $O_p$ bound
alone would not establish the same expected-length claim.

For balanced source sizes, $L_d/(n_s\kappa)=o(1)$, bounded $B_t$ and
$\kappa$ bounded away from zero, the preceding design bounds give
\[
 \E\sqrt{U_D}=O(n_s^{-1/2}K^{-1/4}).
\]
The same order holds for local dispersion $D=O(n_s^{-1}K^{-1/2})$,
including the studied weak-shift sequence. Positive valid-source proportions
are required; contrast certificates additionally require a positive guaranteed
intersection proportion. Arm certificates remain available without that
intersection.

Conditional intersection or shorter fallback never exceeds the source-band
length when source fits are available. Population enlargement is uniformly
$O(n_t^{-1/2})$ for the bounded range/variance methods. Let
$r_n=n_s^{-1/2}K^{-1/4}+n_t^{-1/2}$. Since reported intervals are clipped to
$[-1,1]$,
\[
 \E|I|\le C r_n+2\Prb(\hbox{bad design event}).
\]
Thus merely showing fallback probability tends to zero does not prove the
expected-width rate. A sufficient stronger requirement is
$q_{n,p,K}+\delta_t=O(r_n)$; using $o(r_n)$ makes the fallback contribution
negligible. This distinguishes unconditional coverage, probability-order width,
and expected width. For fixed $p$, the empirical-Bernstein composition method
permits replacing the target term by
$\sigma_\tau n_t^{-1/2}+n_t^{-1}$, using the expectation argument in
Section~\ref{sec:target-protection}; its fallback probability must then be
controlled on that smaller total-width scale.

For approximate models, retaining these rates also requires the valid-source
mean-bias allowance to be negligible relative to $r_n$, the variance-centering
allowance $a_KC_B$ to be small on the $n_s^{-1}K^{-1/2}$ dispersion scale, and
the added concentration constant to have the corresponding order. A fixed
nonzero approximation envelope need not provide improving width as $K$ grows.

\subsection{Design-only removal as a separate possible policy}

If $m$ sources fail an outcome-independent design rule, retaining the others
leaves $K'=K-m$ sources and at least $g_a'=\max(0,g_a-m)$ valid sources per
arm. Conditional inference could be rerun with those conservative counts when
positive. This follows by removing at most $m$ members of each unknown valid
set; it does not require knowing which discarded sites were valid.

The current main estimator still requires all source designs. The stability
study records only whether this count-adjusted alternative would remain
\emph{eligible}; it does not claim that the alternative was fitted or that its
intervals are shorter. It is a concrete way to investigate the fixed-$n_s$,
large-$K$ boundary after the present all-designs-required analysis.
